\slajdid{08-s038}
\begin{frame}[label=08-s038]
\gusttekst
Jasno je da će postojati vreme uspona i pada i na izlaznom signalu i principski su različiti od vremena uspona i pada na ulaznom signala. Međutim samo merenjem vremena usponske i silazne ivice na izlaznom signalu neće nam dati informaciju koliko je uticaj samog kola. Isto tako i neidealna usponska ili silazna ivica na ulaznom signalu utiče na izmereno kašnjenje kola. Zbog toga se pribegava aproksimativnim formulama\par
Kašnjenje silazne ivice, (u slučaju invertora)
\[t_{pHL}=\sqrt{t_{pHL0}^{2}+\left(\frac{t_r}2\right)^2}\]
$t_{pHL}$ – izmereno\par
$t_r$ – izmereno na ulaznom signalu\par
$t_{pHL0}$ – kašnjenje koje potiče od samog kola, i to računamo\par
Ista, slična formula, se koristi i za kašnjenje uzlazne ivice
\[t_{pLH}=\sqrt{t_{pLH0}^{2}+\left(\frac{t_f}2\right)^2}\]
\end{frame}
